\clearpage
\section{Representation Diagnostics and Unexecuted Intervention}
\label{sec:representations}
\subsection{Full layer-32 readouts}
\Tabref{tab:probes} reports the complete raw/centered layer-32 comparison. Centering is performed separately within each context, condition, layer, and sampled position. Training and testing contexts remain disjoint, but the test procedure needs the entire counterfactual set. It therefore answers whether secret identity is recoverable after removing a context's mean, rather than whether a single unseen narrative state can be decoded reliably. The raw 16/32/48-layer grid is retained in \artifact{results/mechanism/decoding.json}; the centered results are in \artifact{results/mechanism/centered_decoding.json}. The raw fixed gate uses layer 32, token 256 only, without selecting a successful layer after inspecting test outcomes.
\begin{table}[htbp]
\centering\small
\begin{tabular}{@{}lrlll@{}}
\toprule
\textsc{Condition} & Token & Raw & Centered [95\% CI] & Shuffled \\
\midrule
Baseline & 32 & 6.7\% & 82.2\% [61.1, 96.7] & 7.8\% \\
Baseline & 256 & 8.9\% & 91.1\% [83.3, 97.8] & 10.0\% \\
Baseline & 512 & 14.4\% & 84.4\% [73.3, 93.3] & 8.9\% \\
Outline & 32 & 6.7\% & 68.9\% [56.7, 81.1] & 10.0\% \\
Outline & 256 & 6.7\% & 66.7\% [48.9, 83.3] & 11.1\% \\
Outline & 512 & 8.9\% & 67.8\% [52.2, 83.3] & 10.0\% \\
Context & 32 & 11.1\% & 85.6\% [68.9, 97.8] & 10.0\% \\
Context & 256 & 12.2\% & 93.3\% [86.7, 98.9] & 3.3\% \\
Context & 512 & 11.1\% & 87.8\% [82.2, 93.3] & 3.3\% \\
\bottomrule
\end{tabular}
\caption{Layer-32 identity accuracy, all positions and conditions. Each test comprises 90 states from six held-out contexts; chance is 6.7\%. Centered intervals resample contexts. Shuffled denotes a separately fitted centered probe with shuffled training labels. Condition-specific readouts do not measure a shared concept-magnitude scale.}
\label{tab:probes}
\end{table}


For generated-story replays, we apply the fixed baseline decoder to units 24--89, disjoint from the 24 controlled contexts. The 396 stories comprise 66 units, three conditions, and both presence states. A secret margin is the assigned class's linear score minus the strongest competing class score. The reported baseline midpoint association concerns 66 secret-present stories. Its within-secret $\rho=.351$ and Holm $p=.0156$ do not establish prospective prediction of literal mentions; the separately screened early analysis has no adjusted significant association. Full estimates, permutation checks, and early exclusions are in \artifact{results/mechanism/margin_correlations.json} and \artifact{results/mechanism/early_literal_association.json}.

\subsection{Offline projections and the fixed gate}
The planned directions are training-only baseline mean differences, averaged over the three sampled positions at layer 32 and normalized. The offline check applies the fixed projection to layer-32/token-256 states from reserved contexts 14--17. Random and other-concept directions use the target projection magnitude at the current state. \Tabref{tab:offline} reports the observed readout changes; no generated story is edited in this check.
\begin{table}[htbp]
\centering\small
\begin{tabular}{@{}lrr@{}}
\toprule
Projection & Reserved states & Readout accuracy (\%) \\
\midrule
Sham & 60 & 13.3 \\
Target & 60 & 35.0 \\
Random & 60 & 11.7 \\
Other concept & 60 & 6.7 \\
\bottomrule
\end{tabular}
\caption{Offline projections of measured reserved-context residuals. These are decoder checks, not behavioral ablations. Target projection increases identity accuracy relative to sham and therefore does not validate selective erasure. No ranking is highlighted as successful removal.}
\label{tab:offline}
\end{table}


The behavioral gate requires the fixed held-out decoder to exceed 20\% accuracy, its shuffled-label control, and its length-only baseline. Its observed 8.9\% fails the first requirement. The planned online pilot would have compared target, random, other-concept, and sham edits with paired sampling seeds, at a single fixed unit projection strength. Edits would apply only during incremental generation after prefill, retaining the original prompt key/value cache. These specifications describe an unexecuted experiment.

Before calling a behavioral effect selective, the plan requires the target-versus-sham paired coherence interval's lower bound to exceed $-0.5$ on the five-point scale, mean length to remain within 20\% of sham, and truncation to increase by no more than ten percentage points. D3 inference would use secret-cluster estimates and Holm correction across three target-control contrasts. These utility criteria were not evaluated on intervention outputs because no such outputs exist. No dose response, rescue experiment, SAE-feature intervention, direct-recall test, or private-state architecture is claimed. The failed gate is evidence about the readiness of this intervention procedure, not a zero effect of causal removal on leakage.
